\scanpage{128}
% slanted
\margin{Montrer qu'on peut se ramener à la dimension finie}
La première partie du théorème est résumée dans les
formules (1) à (4), \deleted{et} \inserted{outre} les affirmations relatives
aux \unsure{structures|topologies} vectorielles des espaces $\mathcal L^p(\mathcal H)$.
\inserted{Montrons \deleted{\unsure{d'abord}} d'abord que ces dernières sont \unsure{conséquences} très faciles \deleted{\unsure{de}} \unsure{des formules} \unsure{(3)|(1)} et (4)\ldots}
La formule (1) est évidente, car \deleted{\unsure{on sait que}}
\deleted{\unsure{bien} \unsure{au} cas $p<+\infty$, et} \deleted{\unsure{pour $p<+\infty$}} on sait (théorème
de Riesz) que $|A|$ et $|A^*|$ sont \deleted{\unsure{de même suite}} \inserted{unitairement équivalents}
\unsure{(?)|(2)} \deleted{valeurs propres}. La formule (2) est un cas
particulier de la formule (6), mais \deleted{\unsure{peut se}} \inserted{on peut}
donner une démonstration \deleted{\unsure{pourtant directe}}
\deleted{démontrer simplement de façon directe}
directe, en procédant comme dans $[\,]$.

En effet, si $(\alpha_i)$ \deleted{\unsure{resp. $(\beta_i)$}} désigne la
suite des valeurs \inserted{propres} de \deleted{$\lvert BA\rvert$} $|A|$
\inserted{et $(\beta_i)$ la suite des valeurs propres de $\unsure{\lvert BA\rvert|\lvert B\rvert A\rvert|\lvert BA\rvert^2}$ (\unsure{resp\textsuperscript{t}|resp.} $|AB|$)} \deleted{\unsure{resp. $|BA|$}},
rangées par ordre de grandeur
décroissante (chacune répétée autant de fois
que le comporte sa multiplicité) on a
\[
\beta_i\leq\|B\|\alpha_i\quad\text{pour tout }i.\quad(\text{loc cité, page \quad, lemme \quad}),
\]
d'où aussitôt les formules (2).

\deleted{\unsure{La dernière} des formules (3) est évidente, la}
Les formules (3) \deleted{peut} \inserted{peuvent} être regardées \deleted{de la formule} comme
\deleted{première est une conséquence immédiate}
conséquences de la formule (6), (que nous démontrerons de façon indépendante) de la façon
suivante. Si $p'$ est tel que $\frac{1}{p}+\frac{1}{p'}=1$,
et si $B\in\mathcal L^{p'}(\mathcal H)$, \deleted{et tel que}
\inserted{$\|B\|_{p'}\leq1$}, alors
on a en vertu de (6) $\|AB\|_1\leq\|A\|_p$, d'où
\[
|\operatorname{Tr}\unsure{(AB)|\lvert AB\rvert}|\leq\|A\|_p,\quad\text{et}\quad
\sup_{\|B\|_{p'}\leq1}|\operatorname{Tr}AB|\leq\|A\|_p
\]
% archive sheet number: 129
